\par\Needspace{6\baselineskip}\medskip
\hypertarget{native-130-article096}{}\pdfbookmark[2]{96.}{article-096}
\begin{center}96.\end{center}
Die im vorhergehenden Artikel angegebenen Gleichungen sind so elegant, dass sie nichts weiter zu wünschen übrig zu lassen scheinen. Dennoch lassen sich einige andere Formeln herleiten, mit denen man die Bahnelemente noch eleganter und bequemer bestimmt. Die Herleitung dieser Formeln ist allerdings etwas weniger naheliegend.

Wir greifen die folgenden Gleichungen aus Art. 8 wieder auf und geben ihnen der Einfachheit halber neue Nummern.
\[
\mathrm{I.}\qquad \sin\tfrac12 v\mathpunct{.}\BareRoot{\dfrac{r}{a}}=\sin\tfrac12 E\mathpunct{.}\BareRoot{(1+e)}
\]
\SourceLeaf{131}
\[
\begin{aligned}
\mathrm{II.}\qquad \cos\tfrac12 v\mathpunct{.}\BareRoot{\dfrac{r}{a}}&=\cos\tfrac12 E\mathpunct{.}\BareRoot{(1-e)}\\[6pt]
\mathrm{III.}\qquad \sin\tfrac12 v'\mathpunct{.}\BareRoot{\dfrac{r'}{a}}&=\sin\tfrac12 E'\mathpunct{.}\BareRoot{(1+e)}\\[6pt]
\mathrm{IV.}\qquad \cos\tfrac12 v'\mathpunct{.}\BareRoot{\dfrac{r'}{a}}&=\cos\tfrac12 E'\mathpunct{.}\BareRoot{(1-e)}.
\end{aligned}
\]
Wir multiplizieren I mit $\sin\frac12(F+g)$ und II mit $\cos\frac12(F+g)$. Addiert man die Produkte, so erhält man
\[
\begin{aligned}
\cos\tfrac12(f+g)\mathpunct{.}\BareRoot{\dfrac{r}{a}}&=\sin\tfrac12 E\sin\tfrac12(F+g)\mathpunct{.}\BareRoot{(1+e)}\\[4pt]
&\quad+\cos\tfrac12 E\cos\tfrac12(F+g)\mathpunct{.}\BareRoot{(1-e)}
\end{aligned}
\]
oder wegen $\BareRoot{(1+e)}=\cos\frac12\varphi+\sin\frac12\varphi$ und $\BareRoot{(1-e)}=\cos\frac12\varphi-\sin\frac12\varphi$
\[
\begin{aligned}
\cos\tfrac12(f+g)\mathpunct{.}\BareRoot{\dfrac{r}{a}}&=\cos\tfrac12\varphi\cos(\tfrac12 F-\tfrac12 G+g)\\[4pt]
&\quad-\sin\tfrac12\varphi\cos\tfrac12(F+G).
\end{aligned}
\]
Ganz entsprechend multipliziert man III mit $\sin\frac12(F-g)$ und IV mit $\cos\frac12(F-g)$ und addiert die Produkte. Das ergibt
\[
\begin{aligned}
\cos\tfrac12(f+g)\mathpunct{.}\BareRoot{\dfrac{r'}{a}}&=\cos\tfrac12\varphi\cos(\tfrac12 F-\tfrac12 G-g)\\[4pt]
&\quad-\sin\tfrac12\varphi\cos\tfrac12(F+G).
\end{aligned}
\]
Subtrahiert man von dieser Gleichung die vorhergehende, so erhält man
\[
\cos\tfrac12(f+g)\mathpunct{.}\left(\BareRoot{\dfrac{r'}{a}}-\BareRoot{\dfrac{r}{a}}\right)=2\cos\tfrac12\varphi\sin g\sin\tfrac12(F-G)
\]
oder nach Einführung des Hilfswinkels $\omega$
\[
\begin{aligned}
\mathrm{27)}\qquad &\cos\tfrac12(f+g)\operatorname{tang}2\omega\\[4pt]
&=\sin\tfrac12(F-G)\cos\tfrac12\varphi\sin g\mathpunct{.}\BareRoot[4]{\dfrac{aa}{rr'}}\mathpunct{\cdot}
\end{aligned}
\]
Durch ganz ähnliche Umformungen, deren Ausführung wir dem kundigen Leser überlassen, findet man
\[
\begin{aligned}
\mathrm{28)}\qquad &\dfrac{\sin\frac12(f+g)}{\cos 2\omega}\\[4pt]
&=\cos\tfrac12(F-G)\cos\tfrac12\varphi\sin g\mathpunct{.}\BareRoot[4]{\dfrac{aa}{rr'}}\\[8pt]
\mathrm{29)}\qquad &\cos\tfrac12(f-g)\operatorname{tang}2\omega\\[4pt]
&=\sin\tfrac12(F+G)\sin\tfrac12\varphi\sin g\mathpunct{.}\BareRoot[4]{\dfrac{aa}{rr'}}\\[8pt]
\mathrm{30)}\qquad &\dfrac{\sin\frac12(f-g)}{\cos 2\omega}\\[4pt]
&=\cos\tfrac12(F+G)\sin\tfrac12\varphi\sin g\mathpunct{.}\BareRoot[4]{\dfrac{aa}{rr'}}\mathpunct{\cdot}
\end{aligned}
\]
Da die linken Seiten dieser vier Gleichungen bekannte Größen sind, bestimmt man aus 27 und 28 die Größen $\frac12(F-G)$ und $\cos\frac12\varphi\sin g\mathpunct{.}\BareRoot[4]{\dfrac{aa}{rr'}}=P$ und ebenso aus 29 und 30 die Größen $\frac12(F+G)$ und $\sin\frac12\varphi\sin g\mathpunct{.}\BareRoot[4]{\dfrac{aa}{rr'}}=Q$. Die Mehrdeutigkeit bei der \SourceLeaf{132}Bestimmung der Winkel $\frac12(F-G)$ und $\frac12(F+G)$ ist so aufzulösen, dass $P$ und $Q$ dasselbe Vorzeichen wie $\sin g$ haben. Aus $P$ und $Q$ leitet man dann $\frac12\varphi$ und $\sin g\mathpunct{.}\BareRoot[4]{\dfrac{aa}{rr'}}=R$ ab. Aus $R$ lassen sich $a=\dfrac{RR\BareRoot{rr'}}{\sin g^2}\mathpunct{\raisebox{.45ex}{,}}$ und $p=\dfrac{\sin f^2\BareRoot{rr'}}{RR}\mathpunct{\raisebox{.45ex}{,}}$ ableiten, sofern man diese Größe, die sich zu $=\pm\BareRoot{\{2(l+\sin\frac12 g^2)\cos f\}}$, $=\pm\BareRoot{\{-2(L-\sin\frac12 g^2)\cos f\}}$ ergeben muss, nicht lieber allein zur Kontrolle der Rechnung verwendet. In diesem Fall bestimmt man $a$ und $p$ am bequemsten nach den Formeln
\[
b=\dfrac{\sin f\mathpunct{.}\BareRoot{rr'}}{\sin g}\mathpunct{\raisebox{.45ex}{,}}\qquad a=\dfrac{b}{\cos\varphi}\mathpunct{\raisebox{.45ex}{,}}\qquad p=b\cos\varphi.
\]
Zur Kontrolle der Rechnung kann man nach Belieben auch mehrere Gleichungen aus Art. 88 und 95 heranziehen. Diesen fügen wir noch die folgenden hinzu.
\[
\begin{aligned}
\dfrac{2\operatorname{tang}2\omega}{\cos 2\omega}\BareRoot{\dfrac{rr'}{aa}}&=e\sin G\sin g\\[6pt]
\dfrac{2\operatorname{tang}2\omega}{\cos 2\omega}\BareRoot{\dfrac{pp}{rr'}}&=e\sin F\sin f\\[6pt]
\dfrac{2\operatorname{tang}2\omega}{\cos 2\omega}&=\operatorname{tang}\varphi\sin G\sin f=\operatorname{tang}\varphi\sin F\sin g.
\end{aligned}
\]
Schließlich erhält man die mittlere Bewegung und die Epoche der mittleren Anomalie ebenso wie im vorhergehenden Artikel.
